\documentclass[../../main.tex]{subfiles}
\graphicspath{{../../image/}{../../figure/}}
\begin{document}

\section{古典解的存在性与破裂}\label{section:ch07-古典解破裂}

当方程组\eqref{align:ch07-典范守恒律}配有凸熵时，\cref{theorem:ch05-凸熵局部存在}
保证存在唯一的、局部定义的古典解，其初值 $U_0$ 位于 Sobolev 空间 $H_2$ 中。然而在一维空间中，存在一个更尖锐的存在性理论，它一般适用于拟线性双曲组（可能是也可能不是守恒律），且不依赖于熵的存在：

\begin{theorem}\label{theorem:ch07-局部存在与破裂}
假定\eqref{align:ch07-典范守恒律}在 $\mathscr{O}$ 上严格双曲。对任一初值 $U_0\in C^1(-\infty,\infty)$（取值于 $\mathscr{O}$ 的紧子集且导数有界），存在一个唯一的 $C^1$ 解，在带形 $(-\infty,\infty)\times[0,T_\infty)$ 上以 $U_0$ 为初值，其中 $0<T_\infty\leqslant\infty$，取值于 $\mathscr{O}$。进而，若 $T_\infty<\infty$，则当 $t\uparrow T_\infty$ 时，$\|\partial_x U(\cdot,t)\|_{L^\infty}\to\infty$ 且/或 $U(\cdot,t)$ 的值域从 $\mathscr{O}$ 的每个紧子集逸出。
\end{theorem}

上述定理的证明可在\cref{section:ch07-注记}
所引文献中找到，它依赖于沿特征线监测 $U$ 及其导数的演化而得到的 $U$ 与 $\partial_x U$ 的逐点界。这类估计将在下文建立，但它们将用于确立有限时间内古典解的破裂。

我们已经遇到过若干古典解破裂的例子，特别是第 6.1 节中标量守恒律的破裂，以及第 7.7 节中弱前沿的破裂。破裂也在压缩简单波存在时发生。事实上，如第 7.6 节所示，一个 $i$-简单波解 $U$ 通过取常微分方程\eqref{align:ch07-稀疏波曲线ODE}的（光滑）解 $V_i$ 与 Burgers 方程\eqref{align:ch07-斜率Burgers}的古典解 $\xi$ 的复合\eqref{align:ch07-简单波参数化}而得到。当\eqref{align:ch07-斜率Burgers}的解破裂时，该 $i$-简单波也破裂。上述例子都只涉及单一特征族。本节的目的是证明：在真正非线性性存在时，来自不同特征族的波之间的相互作用不能阻止光滑解的破裂。

\eqref{align:ch07-典范守恒律}在 $(-\infty,\infty)\times[0,T)$ 上的任一古典 $C^2$ 解 $U$ 可写为
\begin{align}
\begin{cases}
\partial_x U=\displaystyle\sum_{j=1}^n a_j R_j(U)\\
\partial_t U=-\displaystyle\sum_{j=1}^n a_j\lambda_j(U)R_j(U),
\end{cases}
\label{align:ch07-解特征分解}
\end{align}
其中
\begin{align}
a_j=L_j(U)\partial_x U,\qquad j=1,\cdots,n.
\label{align:ch07-aj定义}
\end{align}
由\eqref{align:ch07-简单波分解}，可把\eqref{align:ch07-解特征分解}解释为 $U$ 分解为简单波（每个特征族一个），其强度分别为 $a_1,\cdots,a_n$。我们的目的是沿与 $U$ 关联的 $i$-特征线研究 $a_i$ 的演化。我们令
\begin{align}
\frac{\mathrm{d}}{\mathrm{d}t}=\partial_t+\lambda_i\partial_x
\label{align:ch07-特征方向导数}
\end{align}
表示 $i$-特征方向上的微分。合并\eqref{align:ch07-aj定义}与\eqref{align:ch07-解特征分解}，得到
\begin{align}
\partial_t a_i&=L_i\partial_t\partial_x U+\partial_x U^\top\mathrm{D}L_i^\top\partial_t U\nonumber\\
&=\partial_x(L_i\partial_t U)-\partial_t U^\top\mathrm{D}L_i^\top\partial_x U+\partial_x U^\top\mathrm{D}L_i^\top\partial_t U\nonumber\\
&=\partial_x(L_i\partial_t U)+\sum_{j,k=1}^n(\lambda_j-\lambda_k)R_j^\top\mathrm{D}L_i^\top R_k a_j a_k,
\label{align:ch07-dtai}
\end{align}
\begin{align}
\lambda_i\partial_x a_i&=\partial_x(\lambda_i L_i\partial_x U)-(\mathrm{D}\lambda_i\partial_x U)(L_i\partial_x U)\nonumber\\
&=\partial_x(\lambda_i L_i\partial_x U)-\sum_{j,k=1}^n(\mathrm{D}\lambda_i R_j)\delta_{ik}a_j a_k,
\label{align:ch07-dxai}
\end{align}
其中 $\delta_{ik}$ 是 Kronecker delta。由\eqref{align:ch07-典范守恒律}，$L_i\partial_t U+\lambda_i L_i\partial_x U=0$。此外，由\eqref{align:ch07-左右向量规范化}，$R_j^\top\mathrm{D}L_i^\top R_k=-L_i\mathrm{D}R_j R_k$。因而合并\eqref{align:ch07-特征方向导数}、\eqref{align:ch07-dtai}、\eqref{align:ch07-dxai}并对称化，我们结论
\begin{align}
\frac{\mathrm{d}a_i}{\mathrm{d}t}=\sum_{j,k=1}^n\gamma_{ijk}a_j a_k
\label{align:ch07-riccati}
\end{align}
其中
\begin{align}
\gamma_{ijk}=-\tfrac12(\lambda_j-\lambda_k)L_i[R_j,R_k]-(\mathrm{D}\lambda_i R_j)\delta_{ik},
\label{align:ch07-gamma}
\end{align}
这里 $[R_j,R_k]$ 表示 Lie 括号\eqref{align:ch07-Lie括号定义}。特别地，注意
\begin{align}
\gamma_{iii}=-\mathrm{D}\lambda_i R_i,
\label{align:ch07-gammaiii}
\end{align}
\begin{align}
\gamma_{ijj}=0,\qquad j\neq i.
\label{align:ch07-gammaijj}
\end{align}

显然，在通过\eqref{align:ch07-riccati}展示 $a_i$ 爆破的任一论证中，系数 $\gamma_{iii}$ 将起关键作用。由\eqref{align:ch07-gammaiii}，当 $i$-特征族真正非线性时 $\gamma_{iii}$ 永不消失，而当 $i$-特征族线性退化时 $\gamma_{iii}$ 恒为零。

为获得洞见，我们首先考虑 $U$ 恰为 $i$-简单波的情形，即 $a_i\neq 0$ 而 $j\neq i$ 时 $a_j=0$。在此情形下，\eqref{align:ch07-riccati}退化为
\begin{align}
\frac{\mathrm{d}a_i}{\mathrm{d}t}=\gamma_{iii}a_i^2.
\label{align:ch07-riccati-单波}
\end{align}
进而，由于 $U$ 沿特征线为常数，\eqref{align:ch07-riccati-单波}中的 $\gamma_{iii}$ 是常数。当 $\gamma_{iii}\neq 0$ 且 $a_i$ 与 $\gamma_{iii}$ 同号时，\eqref{align:ch07-riccati-单波}在有限时间内诱导 $a_i$ 的爆破。

接下来我们考虑值得注意的特殊情形：方程组\eqref{align:ch07-典范守恒律}配有 Riemann 不变量坐标系 $(w_1,\cdots,w_n)$。在此情形下 $L_j=\mathrm{D}w_j$，于是由\eqref{align:ch07-aj定义}，
\begin{align}
a_j=\partial_x w_j.
\label{align:ch07-ajRiemann}
\end{align}
进而，由\eqref{align:ch07-gamma}、\eqref{align:ch07-Lie括号零}与\eqref{align:ch07-链式法则}，\eqref{align:ch07-riccati}退化为
\begin{align}
\frac{\mathrm{d}a_i}{\mathrm{d}t}=-\sum_{j=1}^n\frac{\partial\lambda_i}{\partial w_j}a_i a_j.
\label{align:ch07-riccati-Riemann}
\end{align}

我们为\eqref{align:ch07-riccati-Riemann}寻求一个积分因子。若 $\phi$ 是 $U$ 的任一光滑标量函数，由\eqref{align:ch07-解特征分解}我们得到：
\begin{align}
\frac{\mathrm{d}\phi}{\mathrm{d}t}=\mathrm{D}\phi(\partial_t U+\lambda_i\partial_x U)=\sum_{j\neq i}(\lambda_i-\lambda_j)(\mathrm{D}\phi R_j)a_j=\sum_{j\neq i}(\lambda_i-\lambda_j)\frac{\partial\phi}{\partial w_j}a_j.
\label{align:ch07-dphi}
\end{align}
合并\eqref{align:ch07-riccati-Riemann}与\eqref{align:ch07-dphi}，得到
\begin{align}
\frac{\mathrm{d}}{\mathrm{d}t}(e^\phi a_i)=-e^\phi\frac{\partial\lambda_i}{\partial w_i}a_i^2-\sum_{j\neq i}e^\phi\left[\frac{\partial\lambda_i}{\partial w_j}-(\lambda_i-\lambda_j)\frac{\partial\phi}{\partial w_j}\right]a_i a_j.
\label{align:ch07-积分因子}
\end{align}
由\eqref{align:ch07-gjk定义}与\eqref{align:ch07-g对称}，存在 $\phi$ 满足
\begin{align}
\frac{\partial\phi}{\partial w_j}=\frac{1}{\lambda_i-\lambda_j}\frac{\partial\lambda_i}{\partial w_j},\qquad j=1,\cdots,i-1,i+1,\cdots,n.
\label{align:ch07-phi方程}
\end{align}
对这个 $\phi$，\eqref{align:ch07-积分因子}退化为
\begin{align}
\frac{\mathrm{d}}{\mathrm{d}t}(e^\phi a_i)=-e^{-\phi}\frac{\partial\lambda_i}{\partial w_i}(e^\phi a_i)^2.
\label{align:ch07-最终Riccati}
\end{align}

当 $i$-特征族真正非线性时，$\partial\lambda_i/\partial w_i\neq 0$。只要 $e^{-\phi}\partial\lambda_i/\partial w_i$ 在解的值域上一致地远离零，\eqref{align:ch07-最终Riccati}就将沿任一从 $x$ 轴上使 $a_i$ 与 $\partial\lambda_i/\partial w_i$ 异号的点 $\bar{x}$ 发出的特征线，在有限时间内诱导 $a_i$ 的爆破。$e^{-\phi}\partial\lambda_i/\partial w_i$ 的一致有界性得以维持，因为由\hyperref[theorem:ch07-Riemann不变量沿特征线常值]{定理 7.3.4}，Riemann 不变量状态空间中任一古典解的值域与其初值的值域重合。在相反情形，即 $i$-特征族线性退化而 $\partial\lambda_i/\partial w_i$ 恒为零时，\eqref{align:ch07-最终Riccati}蕴含只要解存在 $\partial_x w_i$ 就保持有界。我们由此确立：

\begin{theorem}\label{theorem:ch07-Riemann系破裂}
假定\eqref{align:ch07-典范守恒律}配有 Riemann 不变量坐标系 $(w_1,\cdots,w_n)$。假设 $i$-特征族真正非线性。那么任一初值 $U_0$ 有界、且 $\mathrm{d}w_i(U_0)/\mathrm{d}x$ 在某点 $\bar{x}\in(-\infty,\infty)$ 处与 $\partial\lambda_i/\partial w_i$ 异号的古典解 $U$，在有限时间内破裂。
\end{theorem}

\begin{theorem}\label{theorem:ch07-线性退化全局存在}
假定严格双曲方程组\eqref{align:ch07-典范守恒律}是线性退化的，并在区域 $\mathscr{O}=\{U\in\mathbb{R}^n:a_i<w_i<b_i,\ i=1,\cdots,n\}$ 上配有 Riemann 不变量坐标系 $(w_1,\cdots,w_n)$。对任一取值于 $\mathscr{O}$ 的初值 $U_0\in C^1(-\infty,\infty)$，上半平面 $(-\infty,\infty)\times[0,\infty)$ 上的 Cauchy 问题存在唯一的 $C^1$ 解 $U$。
\end{theorem}

我们现在回到一般情形。当 $i$-特征场真正非线性时，从而由\eqref{align:ch07-gammaiii}有 $\gamma_{iii}\neq 0$，\eqref{align:ch07-riccati}中的项 $\gamma_{iii}a_i^2$ 将具有去稳定效应。任何关于它可能被\eqref{align:ch07-riccati}中其余项（它们解释与其他特征场的相互作用效应）所抵消的预期，都不太可能实现，至少当初值在有界区间外为常数时如此，理由如下。方程\eqref{align:ch07-gammaijj}排除了其余特征场自相互作用的可能性：除 $\gamma_{iii}a_i^2$ 外的所有相互作用都涉及两个不同的特征族。现在，当初值在有界区间外为常数时，相互作用最终变得不显著，因为不同特征族的波以不同速度传播并最终分离。因而在长时间中，项 $\gamma_{iii}a_i^2$ 成为主导因子，并在有限时间内把 $a_i$ 驱至无穷。上述启发式论证可以被形式化，并导出下述定理：

\begin{theorem}\label{theorem:ch07-小扰动破裂}
假定\eqref{align:ch07-典范守恒律}是真正非线性的严格双曲守恒律组。考虑初值 $U_0\in C^2(-\infty,\infty)$，使得当 $x\leqslant a$ 与 $x\geqslant b$ 时 $U_0(x)=\bar{U}$（常状态）。若 $(b-a)^2\max|U_0''(x)|$ 是充分小的正数，则初值问题的古典解在有限时间内破裂。
\end{theorem}

在\cref{section:ch07-弱解}
所引文献中，读者将找到上述定理的（冗长而技术性的）证明，以及在更弱假设下古典解破裂的类似结果，即当放松方程组严格双曲性的要求、仅某些特征族真正非线性时。还存在\hyperref[theorem:ch07-线性退化全局存在]{定理 7.8.3}的推广，其中当特征族关于某常状态 $\bar{U}$ 仅是\emph{弱线性退化}时（即对 $i=1,\cdots,n$，$\mathrm{D}\lambda_i R_i$ 只需沿从 $\bar{U}$ 发出的 $i$-稀疏波曲线为零），且初值 $U_0\in C^1(-\infty,\infty)$ 充分靠近 $\bar{U}$（其意义为
\begin{align}
\sup\left\{(1+|x|)^{1+\mu}\left(|U_0(x)-\bar{U}|+|U_0'(x)|\right)\right\}
\label{align:ch07-小初值}
\end{align}
对某 $\mu>0$ 充分小）时，建立了 $C^1$ 解的全局存在性。

关于初值--边值问题也有大量研究。为与第 5.6 节的讨论相比较，我们考虑严格双曲方程组\eqref{align:ch07-典范守恒律}在象限 $\{(x,t):x>0,\ t>0\}$ 上，假定边界 $x=0$ 是非特征的，即对某 $k=0,\cdots,n$，$\lambda_k(U)<0<\lambda_{k+1}(U)$，其中 $\lambda_0(U)=-\infty$，$\lambda_{n+1}(U)=\infty$。我们施加形式为\eqref{align:ch05-ibvp-boundary}
的边界条件，即 $BF(U(0,t))=0$，其中 $B$ 是 $n\times n$ 矩阵，使得对流形 $\mathscr{M}=\{U:BF(U)=0\}$ 上任一 $U$，$\mathbb{R}^n$ 是 $\mathrm{BD}F(U)$ 的核与 incoming 特征族 $\{R_{k+1}(U),\cdots,R_n(U)\}$ 所张成子空间的直和。我们还指定与边界条件相容的初值 $U_0\in C^1([0,\infty);\mathbb{R}^n)$，即 $BF(U_0(0))=0$ 且 $B[\mathrm{D}F(U_0(0))]^2U_0'(0)=0$。在上述假设下，可以证明：(a) 古典解的局部存在性；(b) 当（至少某些）incoming 特征族真正非线性时古典解在有限时间内破裂；(c) 当所有 incoming 特征族（至少）弱线性退化且初值的上确界\eqref{align:ch07-小初值}充分小时古典解的全局存在性。

\end{document}
